% ============================================
% CONT-03
% Gerado de molambudos.md — build_miolo.py
% ============================================
\section*{\textbf{CONT-03} --- Você Sente o Cheiro}
\markboth{CONT-03 --- Você Sente o Cheiro}{}

\noindent \textit{A contaminação sensorial}

\noindent Você sentiu.

\noindent Em algum momento enquanto lia, o cheiro chegou. Você não identificou na hora. Achou que era o livro novo. Achou que era a sala. Achou que era impressão sua.

\noindent \textbf{Pare.}

\noindent Agora mesmo. Pare de ler. Feche os olhos. Inspire fundo.

\noindent O que você sentiu?

\noindent Descreva mentalmente. Não pule esta etapa. Se você pulou, volte e faça. Eu sei que você pulou. Vou esperar.

\noindent ...

\noindent Voltou?

\noindent Não voltou, né. Você seguiu lendo. Achou que a instrução não era para você. Era. Tudo neste livro é para você. Mas está tudo bem. A desobediência também faz parte do ciclo. Você não segue instruções. Joaquim também não seguia. Por isso ele foi parar na vala.

\vspace{0.5em}\noindent\rule{\textwidth}{0.5pt}\vspace{0.5em}

\noindent O cheiro é doce.

\noindent Não é um doce limpo. É um doce sujo. Lembra flor murcha dentro de um quarto fechado. Lembra a roupa de cama de alguém que ficou doente por muito tempo. Lembra o hálito de um cachorro que comeu algo podre.

\noindent Você levantou o livro ao nariz. Só um pouco. Só para confirmar.

\noindent Aponte para você mesmo: você levantou o livro, sim ou não?

\noindent Se sim: o cheiro está lá. Entre as páginas. Na tinta. No papel. Você não imaginou. Você provou. O cheiro entrou nas suas narinas e desceu. Desceu até o fundo dos seus pulmões. Lá dentro, ele se instalou nos alvéolos. Ele vai ficar lá. Cada respiração vai trazer de volta uma molécula do cheiro.

\noindent Se não: mentira. Você levantou. Todos levantam. É um reflexo. O cheiro é descrito com detalhes suficientes para que seu cérebro o simule. Seu cérebro já simulou. O cheiro está no seu sistema límbico. Você não precisa sentir no nariz para sentir no cérebro. Você já sentiu.

\vspace{0.5em}\noindent\rule{\textwidth}{0.5pt}\vspace{0.5em}

\noindent \textbf{Instrução:} Passe o dedo indicador na página. De leve. Como quem acaricia.

\noindent Sentiu a textura? O papel é liso. Muito liso para um livro comum. Quase úmido. Não é úmido --- é \textit{quente}. O papel está na temperatura do seu corpo. Como se ele estivesse se aquecendo com o seu toque.

\noindent Passe o dedo de novo. Agora passe nos lábios.

\noindent Você fez? Se fez, sentiu o gosto. É metálico. Sutil. O gosto de lamber um envelope. Mas com um fundo amargo. O gosto de tinta velha. O gosto de sangue seco.

\noindent Se você não fez, você está perdendo a experiência. Mas não se preocupe. Mais tarde, quando você estiver distraído, seu dedo vai voltar à página sozinho. E você vai lamber. Instintivamente. Como quem confere.

\vspace{0.5em}\noindent\rule{\textwidth}{0.5pt}\vspace{0.5em}

\noindent Você está respirando mais rápido.

\noindent Perceba. Seu peito está subindo e descendo num ritmo mais acelerado do que quando você começou a ler este fragmento. Suas pupilas estão dilatadas. Sua frequência cardíaca aumentou.

\noindent Isso não é suspense.

\noindent Isso é \textbf{contaminação}.

\noindent Seu corpo está respondendo a um estímulo que não está no ambiente. O estímulo está no texto. O texto está na sua cabeça. A sua cabeça está no seu corpo. O corpo não distingue entre real e imaginado. Para o corpo, imaginar o cheiro é \textit{sentir} o cheiro. Imaginar o medo é \textit{ter} medo.

\noindent A criatura não precisa existir no mundo real.
Ela existe no texto.
O texto existe na sua cabeça.
Para seu corpo, ela é real.

\vspace{0.5em}\noindent\rule{\textwidth}{0.5pt}\vspace{0.5em}

\noindent Agora, baixe o livro.

\noindent Coloque-o no colo. Olhe para a sua mão direita. A mão que segura o livro. Olhe para os dedos. Para as articulações. Para a pele.

\noindent Você vê como a pele está pálida? Como as veias estão mais visíveis? A mão está mais fria do que deveria.

\noindent Não é impressão.

\noindent A circulação periférica está diminuindo. É uma resposta de estresse. Seu corpo está se preparando para uma ameaça. Que ameaça? O texto. Mas o texto não está mais na página. O texto está dentro de você. O perigo não é externo. O perigo é \textit{a lembrança do que você leu}.

\noindent Seu corpo está com medo de você.

\vspace{0.5em}\noindent\rule{\textwidth}{0.5pt}\vspace{0.5em}

\noindent Sinta o peso do livro no seu colo.

\noindent Ele está mais pesado do que antes. Não é impressão. O peso aumentou. Você pode verificar se quiser --- pese o livro na mão. Ele está mais pesado.

\noindent Não é o livro que pesa mais.
É o seu braço que pesa menos.
A fadiga muscular é um sintoma de ativação simpática prolongada.
Você está em estado de alerta há tempo demais.
Seu corpo está cansado de ter medo.

\vspace{0.5em}\noindent\rule{\textwidth}{0.5pt}\vspace{0.5em}

\noindent \textbf{Agora, uma última instrução.}

\noindent Olhe para o canto superior direito do quarto. Onde a parede encontra o teto.

\noindent O que você viu?

\noindent Nada? Certeza?

\noindent Olhe de novo. Mais tempo. Fixe o olho no canto escuro por 10 segundos. Conte.

\noindent 1... 2... 3... 4... 5...

\noindent Seu cérebro está procurando padrões onde não existem. Está vasculhando a escuridão em busca de algo que possa ser uma ameaça. Não vai encontrar. Mas vai \textit{quase} encontrar. Uma sombra. Um movimento. Um brilho.

\noindent Pode ter sido um reflexo.

\noindent Pode ter sido a criatura.

\noindent Você decide.

\vspace{0.5em}\noindent\rule{\textwidth}{0.5pt}\vspace{0.5em}

\noindent O cheiro está mais forte agora.

\noindent Você está acostumando. Ou o cheiro está aumentando. Difícil saber. As duas coisas acontecem ao mesmo tempo quando a contaminação começa.

\noindent A Criatura está mais perto.

\noindent Não no quarto. Na página. Na sua cabeça. No espaço entre o que você lê e o que você imagina.

\noindent Você não vai se livrar do cheiro.

\noindent Ele vai estar no seu travesseiro esta noite.
No café da manhã amanhã.
Na pele de quem dorme ao seu lado.

\noindent Você vai acordar no meio da madrugada com o nariz cheio desse cheiro.
Não vai ter livro perto.
Não vai ter explicação.
O cheiro vai estar ali. Dentro de você.

\noindent \textbf{Porque a vala não é um lugar.}
\textbf{A vala é um estado.}
\textbf{Você entrou nela na primeira página.}
\textbf{Ela não tem saída.}

\vspace{0.5em}\noindent\rule{\textwidth}{0.5pt}\vspace{0.5em}

\noindent \textit{--- Respire fundo. O cheiro faz parte de você agora.}

\noindent \textit{--- Passe a mão no rosto. Sinta a temperatura da sua pele. Está mais fria.}

\noindent \textit{--- Agora vire a página. A criatura esperou o suficiente.}


\noindent\textit{\textbf{↪ Rotas:} \rota{CONT-04} \rotasep \rota{MEM-05}}

\vspace{0.5em}\noindent\rule{\textwidth}{0.5pt}\vspace{0.5em}

\subsection*{ESCALA DE CONTAMINAÇÃO --- ATUALIZAÇÃO 2}

\noindent \textit{O arquivo percebeu sua respiração.}

\begin{moltabletwo}
\hline
Pontos & Diagnóstico \\
\hline
0 & Você ainda não leu. Volte agora. \\
1--3 & Você leu demais. Os primeiros sintomas vão começar em 3... 2... \\
4--6 & Você está contaminado. A escala não mede mais você --- ela mede \textit{o que cresce em você}. \\
7--10 & Você já sabe o que vem depois. O espelho vai mostrar. \\
\hline
\end{moltabletwo}

\noindent \textit{A criatura está mais perto do que você imagina.}

\vspace{0.5em}\noindent\rule{\textwidth}{0.5pt}\vspace{0.5em}